% Generated draft: requires mathematical and source review.
\documentclass[11pt]{article}
\usepackage[margin=25mm]{geometry}
\usepackage{fontspec}
\setmainfont{Latin Modern Roman}
\usepackage{amsmath,amssymb,mathtools,unicode-math}
\setmathfont{Latin Modern Math}
\usepackage{enumitem,needspace,xurl,hyperref}
\setlength{\parindent}{0pt}
\setlength{\parskip}{5pt}
\setlength{\emergencystretch}{3em}
\newcommand{\N}{\mathbb{N}}
\newcommand{\Z}{\mathbb{Z}}
\newcommand{\Q}{\mathbb{Q}}
\newcommand{\R}{\mathbb{R}}
\newcommand{\C}{\mathbb{C}}
\newcommand{\F}{\mathbb{F}}
\newcommand{\A}{\mathbb{A}}
\newcommand{\PP}{\mathbb{P}}
\newcommand{\E}{\mathbb{E}}
\newcommand{\eps}{\varepsilon}
\newcommand{\dd}{\,\mathrm{d}}
\begin{document}
\section*{UnsolvedMath ID: 1200026}
\textbf{Source problem code:} catalogue entry
\subsection{Definitions and mathematical statement}
Let $\Gamma$ be an infinite finitely generated group which is not virtually cyclic, meaning that it has no cyclic subgroup of finite index. Let $S$ be a finite symmetric generating set and let $G=\operatorname{Cay}(\Gamma,S)$ be the corresponding locally finite undirected Cayley graph. Bernoulli $p$-bond percolation retains each edge independently with probability $p\in[0,1]$. Its clusters are the connected components of the retained-edge graph.

Fix a vertex $o$, put $\theta_G(p)=\mathbb P_p(o\text{ lies in an infinite cluster})$, and define
\[
 p_c(G)=\inf\{p\in[0,1]:\theta_G(p)>0\}.
\]
The requested assertion is
\[
 \forall\Gamma\ \forall S\text{ as above},\qquad p_c(G)<1.
\]
Equivalently, there exists $p\in(0,1)$ such that Bernoulli $p$-bond percolation has an infinite cluster almost surely. The latter is the global existence event, rather than the assertion that the chosen vertex lies in an infinite cluster almost surely.

The original statement only says ``infinite Cayley graph''; finite generation and a finite generating set are made explicit here as the standard locally finite setting used by the supplied resolving assessment. No restriction to a particular generating set, growth rate beyond the stated group assumption, or uniqueness of the infinite cluster is added.
\subsection{Short English statement}
Every infinite finitely generated group outside the virtually cyclic class should have a percolation phase with infinite clusters strictly below edge-retention probability one. The supplied assessment labels this proved: it attributes the superlinear-growth implication to the linear-growth characterization of virtually cyclic groups (Justin/Gromov), the quasi-transitive superlinear-growth percolation theorem to Duminil-Copin--Goswami--Raoufi--Severo--Yadin (2020), and a later proof via uniform transience to Easo--Severo--Tassion (2025). These are cited resolution claims, not a new proof or independent source verification.
\subsection{Sources}
\begin{itemize}
\item UnsolvedMath permanent ID 1200026, source code catalogue entry, ``Some Questions --- Question 26.'' Registry curation and metadata: Ulam AI and UnsolvedMath Contributors, CC BY 4.0; updated 2026-08-06.
\item Abert, \emph{Some questions} (2010), Question 26: \url{https://www.renyi.hu/~abert/questions.pdf}. Supplied manual transcription accessed 2026-07-29. Original metadata says current status was not independently verified there; status/release and rights review are NEEDS\_REVIEW.
\item Duminil-Copin--Goswami--Raoufi--Severo--Yadin, \emph{Duke Mathematical Journal} (2020), and Easo--Severo--Tassion, \emph{Forum of Mathematics, Pi} (2025), as cited in the supplied resolving assessment. Justin/Gromov is named there for the linear-growth characterization. No article titles, volumes, pages, or locators are supplied for these references.
\end{itemize}
\end{document}
